\scanpage{68}
\mspage{33}
\uline{Proposition} Le diviseur à l'infini \inserted{$M^\infty_{0,I}$} de $\widehat M_{0,I}$ est une courbe \unsure{stable|semi-stable} \inserted{déployée} sur $\operatorname{Spec}\mathbf Z$, dont le graphe est isomorphe au \deleted{celui du} $1$-squelette du dodécaèdre gauche défini par $I$ (muni d'une orientation convenable).
\margin{\uline{NB} $\widehat M_{0,I}\longrightarrow$ un vrai schéma projectif et lisse \unsure{et tant} --- le groupoïde qu'il représente est en effet rigide. Il en \unsure{va} pour tous les $M_{0,I}\simeq M^!_{0,\nu}$ quelque soit $\nu\in\mathbf N\ldots$}

Bien entendu, les genres des composantes irréductibles sont nuls. Il n'y a pas \inserted{lieu de mettre} des ``pts marqués'' sur \deleted{\unsure{toutes}} les \inserted{$10$} composantes irréductibles de $M^\infty_{0,I}$, qui est de toutes façons une MD-courbe \deleted{stable} à MD-graphe maximal (tous les genres nuls, tous les sommets d'ordre $3$, pas de pieds marqués).

L'isom. de la prop. précédente tient compte des opérations de $\mathfrak S_I^+$ \inserted{$\simeq\operatorname{Aut}\text{ du dodécaèdre gauche}$} sur $M^\infty_{0,I}$, mais pas de celles de $\mathfrak S_I$ \deleted{tout} entier --- il faudrait le comprendre du pt de vue de la géométrie de l'icosaèdre \inserted{gauche} (ou du dodécaèdre \inserted{gauche}\ldots) --- ça doit être lié bien sûr au \uline{bi-icosaèdre} gauche combinatoire)

Quand on passe à la multiplicité modulaire
\margin{67}
